\scanpage{8}
Enfin, pour exprimer la condition d'holomorphie de la \unsure{fibration|filtration}, il convient d'introduire, \deleted{sur} sur l'ens \inserted{$\unsure{\mathscr{J}|\mathcal{J}|\gamma}$} des $j$ (ou $j'$) possibles, une structure de \uline{variété} \uline{analytique} \uline{complexe} naturelle, et la condition s'exprime par l'\uline{holomorphie} de l'application
\[
s\longmapstoJ(s):\quad S\rightharpoonup\mathscr{J}.
\]

\deleted{En un \unsure{mot}} Donc l'ens des $G$-subordinations \inserted{rigidifiées} base $S$ s'identifie à l'ens. $\operatorname{Mor}(S,\mathscr{J})$\margin{morph. d'espaces analytiques}. Donc pour $S$ variable, on a un foncteur représentable par $\mathscr{J}$.

\uline{NB} $G(\mathbb{R})$ opère de façon évidente sur $\mathscr{J}$ (mod autom intérieurs) ; les orbites sont \uline{ouvertes}. Donc

\deleted{Soit \deleted{\unsure{$G$ un}} $V$ un $G$-module fidèle, et $V_0\subset V$ un $Z$-module de type fini engendrant $V$. Alors $\forall$ $G$-subordination sur un $S$ définit une structure de Hodge (\uline{pas} mod isog) sur $V_0\times S$, Soit $\Gamma\subset G(\mathbb{Q})$ le stabilisateur de $V_0$.}

les comp. connexes de $\mathscr{J}$ sont de la forme $G(\mathbb{R})^{0}/H(\mathbb{R})\cap G(\mathbb{R})^{0}$, où $H$ est le centralisateur d'un $j:\unsure{\mathbb{S}|\mathrm{S}}\longrightarrow G$ dans $G$ ; ou encore $G'(\mathbb{R})^{0}/H'(\mathbb{R})\cap G'(\mathbb{R})^{0}$, $H'$ le centralisateur d'un $j':\mathbb{U}\longrightarrow G'$ dans $G'(\mathbb{R})^{0}$.
% archive sheet number: 7

